\documentclass[../../../include/open-logic-section]{subfiles}

\begin{document}

\olfileid{sth}{card-arithmetic}{ch}

\olsection{સાતત્યક પરિકલ્પના}

અગાઉનું પરિણામ (યોગ્ય રીતે) સૂચવે છે કે અનંત
ગણાંકો $2$ની ઘાત સાથે, એટલે કે
\olref[opps]{lem:SizePowerset2Exp} મુજબ ઘાતગણ
લેવાની ક્રિયા સાથે, ``સીધી રીતે'' વર્તે તેની ખાતરી
હોય, તો ગણાંક ઘાતક્રિયા બહુ \emph{સહેલી} થાય.
પરંતુ અનંત ગણાંકો ઘાતગણ સાથે ખરેખર \emph{એમ}
વર્તે છે એવું માની શકાતું નથી.

આ વાત સમજવા પહેલાં કેટલીક ઉપયોગી સંજ્ઞાઓ
રજૂ કરીએ.

\begin{defn}
$\cardsucc{\cardfont{a}}$ એ $\cardfont{a}$ કરતાં
કડક મોટો લઘુતમ ગણાંક હોય, તો બે અનંત અનુક્રમો
આમ વ્યાખ્યાયિત કરીએ:
\begin{align*}
	\aleph_{0} &\defis \omega & 		
	\beth_{0} &\defis \omega\\
	\aleph_{\alpha \ordplus 1} &\defis \cardsucc{(\aleph_{\alpha})} &
	\beth_{\alpha+1} &\defis \cardexpo{2}{\beth_{\alpha}}\\
	\aleph_{\alpha} &\defis \bigcup_{\beta< \alpha} \aleph_{\beta} &
	\beth_{\alpha} &\defis \bigcup_{\beta < \alpha}\beth_{\beta} & \text{જ્યારે $\alpha$ સીમા ક્રમવાચક સંખ્યા હોય}.
	\end{align*}
\end{defn}

$\cardsucc{\cardfont{a}}$ની વ્યાખ્યા યોગ્ય છે,
કારણ કે \olref[cardinals][classing]{lem:NoLargestCardinal}
મુજબ દરેક ગણાંક $\cardfont{a}$ માટે
$\cardfont{a}$ કરતાં મોટો કોઈ ગણાંક છે, અને
અનંતાતીત અનુમાનથી $\cardfont{a}$ કરતાં મોટા
ગણાંકોમાં \emph{લઘુતમ} ગણાંક મળે છે. આલેફ અને બેથના
અનુક્રમોની બાકીની વ્યાખ્યાઓ અનંતાતીત
પુનરાવર્તનથી મળે છે.
\sourcecorrection{OLMD-140}{મૂળનું છેલ્લું વાક્ય
``$\cardfont{a}$ની બાકીની વ્યાખ્યા'' કહે છે,
જોકે ઉપર વ્યાખ્યાયિત થતી વસ્તુઓ આલેફ અને બેથના
અનુક્રમો છે. અહીં એ વિષયને સ્પષ્ટ કર્યો છે;
ગણિતીય પુનરાવર્તી સૂત્રો બદલ્યાં નથી.}

કૅન્ટરે ``$\aleph$'' સંજ્ઞા રજૂ કરી હતી;
તે \emph{આલેફ} છે, હિબ્રુ મૂળાક્ષરનો પહેલો
અક્ષર અને ``અનંત'' માટેના હિબ્રુ શબ્દનો પહેલો
અક્ષર. પિર્સે ``$\beth$'' સંજ્ઞા રજૂ કરી હતી;
તે \emph{બેથ} છે, હિબ્રુ મૂળાક્ષરનો બીજો
અક્ષર.\footnote{પિર્સે ડિસેમ્બર 1900માં કૅન્ટરને
લખેલા પત્રમાં આ સંજ્ઞા વાપરી હતી. કમનસીબે
ત્યાં તેમણે $\alpha \geq \omega$ માટે
$\beth_\alpha$ અસ્તિત્વમાં નથી એવી ખોટી દલીલ
પણ આપી હતી.} હવે આ સંજ્ઞાઓ અનંત ગણાંકો આપે છે.

\begin{prop}
દરેક ક્રમવાચક સંખ્યા $\alpha$ માટે
$\aleph_\alpha$ અને $\beth_\alpha$ ગણાંકો છે.
\end{prop}

\begin{proof}
બંને પરિણામો સરળ અનંતાતીત અનુમાનથી મળે છે.
$\aleph_0 = \beth_0 = \omega$ ગણાંક છે,
\olref[cardinals][classing]{omegaisacardinal} મુજબ.
$\aleph_\alpha$ અને $\beth_\alpha$ બંને ગણાંકો
છે એમ ધારીએ, તો $\aleph_{\alpha+1}$ અને
$\beth_{\alpha+1}$ વ્યાખ્યાથી જ ગણાંકો છે.
અને ગણાંકોના ગણનો યોગગણ પણ ગણાંક છે,
\olref[cardinals][classing]{unioncardinalscardinal} મુજબ.
\end{proof}
\noindent
વળી દરેક અનંત ગણાંક કોઈ $\aleph$ છે:

\begin{prop}
જો $\cardfont{a}$ અનંત ગણાંક હોય, તો કોઈ અનન્ય
$\gamma$ માટે $\cardfont{a} = \aleph_\gamma$.
\end{prop}

\begin{proof}
અનંત ગણાંકો પર અનંતાતીત અનુમાન કરો. આધાર
કિસ્સામાં લઘુતમ અનંત ગણાંક આલેફ-શૂન્ય છે.
અનુમાન-ધારણા એવી લો કે અનંત
$\cardfont{b} < \cardfont{a}$ હોય,
તો $\cardfont{b} = \aleph_{\gamma_\cardfont{b}}$.
કોઈ $\cardfont{b}$ માટે
$\cardfont{a} = \cardsucc{\cardfont{b}}$ હોય, તો
$\cardfont{a} = \cardsucc{(\aleph_{\gamma_\cardfont{b}})}=
\aleph_{\gamma_\cardfont{b}+1}$. જો $\cardfont{a}$
કોઈ ગણાંકનો અનુગામી ન હોય, તો ગણાંકો ક્રમવાચક
સંખ્યાઓ હોવાથી
$\cardfont{a} = \bigcup_{\omega \leq \cardfont{b} < \cardfont{a}} \cardfont{b} =
\bigcup_{\omega \leq \cardfont{b} < \cardfont{a}}{\aleph_{\gamma_\cardfont{b}}}$;
તેથી $\gamma = \bigcup_{\omega \leq \cardfont{b} <
\cardfont{a}}\gamma_\cardfont{b}$ લઈએ, તો
$\cardfont{a} = \aleph_\gamma$. આ અનુક્રમ કડક
વધતો હોવાથી સૂચક અનન્ય છે.
\sourcecorrection{OLMD-141}{મૂળની અનુમાન-ધારણા
દરેક નાના ગણાંક માટે આલેફ-સૂચક માને છે;
પરંતુ સાન્ત ગણાંકો આલેફ નથી. લઘુતમ અનંત ગણાંકને
અલગ આધાર કિસ્સો લઈને અને સીમા-કિસ્સાના
યોગગણોમાં ફક્ત અનંત ગણાંકો રાખીને દલીલને
સુગઠિત કરી છે. અનન્યતા માટે અનુક્રમની કડક
વૃદ્ધિ પણ સ્પષ્ટ કરી છે.}
\end{proof}

દરેક અનંત ગણાંક કોઈ $\aleph$ છે, તેથી પ્રશ્ન થાય:
શું દરેક અનંત ગણાંક કોઈ~$\beth$ પણ છે? જો
\emph{એવું} હોય, તો અનંત ગણાંકો ઘાતગણ લેવાની
ક્રિયા સાથે ``સીધી રીતે'' વર્તે. ખરેખર નીચેનું
મળે:

\begin{defish}
\emph{સામાન્યીકૃત સાતત્યક પરિકલ્પના} (GCH).
બધા $\alpha$ માટે $\aleph_\alpha  = \beth_\alpha$.
\end{defish}

વળી GCH સાચી હોય, તો ગણાંક ઘાતક્રિયાનાં ઘણાં
સરળીકરણો કરી શકીએ. ખાસ કરીને મોટો ગણાંક અનંત
હોય અને
$0<\cardfont{b} < \cardfont{a}$ હોય ત્યારે
$\cardexpo{\cardfont{a}}{\cardfont{b}}$નું મૂલ્ય
$\cardfont{a}\leq \cardexpo{\cardfont{a}}{\cardfont{b}} \leq
\cardsucc{\cardfont{a}}$ વચ્ચે બંધાય છે તે
બતાવી શકીએ. પછી આ બેમાંથી કઈ શક્યતા મળે છે
તે નક્કી કરતી ચોક્કસ શરતો આપી શકીએ
(એટલે $\cardfont{a} =
\cardexpo{\cardfont{a}}{\cardfont{b}}$ કે
$\cardexpo{\cardfont{a}}{\cardfont{b}} =
\cardsucc{\cardfont{a}}$).\footnote{આ શરત
\emph{સહઅંતિમતા}થી નક્કી થાય છે.}
\sourcecorrection{OLMD-142}{મૂળમાં માત્ર
$\cardfont{b}<\cardfont{a}$ શરત છે, પરંતુ
$\cardfont{b}=0$ હોય ત્યારે
$\cardexpo{\cardfont{a}}{0}=1$ અને અનંત
$\cardfont{a}$ માટે દર્શાવેલી નીચેની ઉચ્ચસીમા
$\cardfont{a}\leq 1$ ખોટી પડે છે. તેથી આ
મર્યાદિત દાવામાં $\cardfont{b}>0$ ઉમેર્યું છે.
મોટો ગણાંક અનંત હોવાની સંદર્ભ-શરત પણ સ્પષ્ટ કરી છે;
સાન્ત $\cardfont{a}=3$, $\cardfont{b}=2$ હોય તો
$3^2=9>4=\cardsucc{3}$ મળે છે.}

પરંતુ GCH એ \emph{પરિકલ્પના} છે,
\emph{પ્રમેય} નથી. ખરેખર \citet{Godel1938}એ
સાબિત કર્યું કે $\ZFC$ સુસંગત હોય, તો
$\ZFC + \text{GCH}$ પણ સુસંગત છે. પણ પછી
ખબર પડી કે $\ZFC$માં $\lnot$GCH પણ એ જ રીતે
ઉમેરી શકીએ. ખરેખર GCHનો સૌથી સરળ બિનતુચ્છ
\emph{દાખલો} જુઓ:

\begin{defish}
\emph{સાતત્યક પરિકલ્પના} (CH).
$\aleph_1 = \beth_1$.
\end{defish}

\citet{Cohen1963}એ સાબિત કર્યું કે $\ZFC$
સુસંગત હોય, તો $\ZFC + \lnot\text{CH}$ પણ
સુસંગત છે. તેથી સાતત્યક પરિકલ્પના $\ZFC$થી
સ્વતંત્ર છે.

આ પરિકલ્પનાને સાતત્યક પરિકલ્પના કહે છે, કારણ કે
``સાતત્યક'' એ વાસ્તવિક રેખા $\Real$નું બીજું નામ છે.
\olref[opps]{continuumis2aleph0} કહે છે કે
$\card{\Real} = \beth_1$. એટલે સાતત્યક
પરિકલ્પનાનું વિધાન છે કે પ્રાકૃતિક સંખ્યાઓના ગણાંક
$\aleph_0 = \beth_0$ અને સાતત્યકના ગણાંક
$\beth_1$ વચ્ચે કોઈ ગણાંક નથી.

$\ZFC$થી (G)CHની \emph{સ્વતંત્રતા} હોય, તો
તેની \emph{સત્યતા} વિશે શું કહેવું? ઘણું કહી શકાય.
ખરેખર ગણસિદ્ધાંતમાં તાજેતરનું ઘણું ફળદાયી કામ
આ પ્રશ્નોની તપાસમાં થયું છે. પરંતુ બે ટૂંકા મુદ્દા
અવશ્ય ભારપૂર્વક કહેવા જેવા છે.

પહેલું: આ ઔપચારિક સ્વતંત્રતાનાં પરિણામોથી GCH
કે CHનું સત્યમૂલ્ય \emph{અનિર્ધારિત} છે એવું
\emph{તરત} અનુસરતું નથી. કદાચ આપણને સ્વાભાવિક
લાગતી વધુ સ્વયંસિદ્ધિઓ ઉમેરવી પડે અને તે પ્રશ્નને
એક રીતે કે બીજી રીતે નક્કી કરે. ગેડલે પોતે આ
પ્રતિભાવ યોગ્ય હોવાનું સૂચવ્યું હતું.

બીજું: $\ZFC$થી CHની સ્વતંત્રતા નિઃસંદેહ
\emph{આશ્ચર્યજનક} છે, પણ શાબ્દિક અર્થમાં
\emph{અવિશ્વસનીય} નથી. મુદ્દો એટલો જ છે કે
$\ZFC$ જે કહે છે તેમાં ગણાંકથી તેના અનુગામી
ગણાંક સુધી પહોંચવાની રીત ફક્ત ઘાતગણ લેવાની
બળજબરીભરી રીત જ હોવી જરૂરી નથી.
 %The operation of taking powersets moves us from one stage of the
 %hierarchy to its successor stage (from $V_{\alpha}$ to
 %$V_{\alpha+1}$).

આ બે અવલોકનો પછી વધુ જાણવા માટે તમારે આ
ચર્ચામાં રસ ધરાવતા વિવિધ તત્ત્વજ્ઞાનીઓ અને
ગણિતજ્ઞો તરફ વળવું પડશે.\footnote{શરૂઆત માટે
\citet[\S15.6]{Potter2004} વાંચી શકો છો.}

\end{document}
